\section{Audit of the available rows and their analytic hypotheses}
\label{sec:row-audit}

\paragraph{Finding R01 (exact source coverage).}
The six attached \texttt{SORTED\_*.md} files are summaries, not the
complete row lists named in them. Their statements are frequently
truncated. They report $1190$ register rows. Subsequent recovery gives
the following exact inventory; ``proof blocks'' means available numbered
source blocks, not independently verified proofs.
\begin{center}
\begin{tabular}{lrrr}
\hline
Run & Reported rows & Recovered rows & Available proof blocks\\
\hline
\texttt{ANABELIAN\_A1} &121&121&121\\
\texttt{ANABELIAN\_A1\_SUCCESSOR}&210&0&0\\
\texttt{ANABELIAN\_SUCCESSOR}&157&157&157\\
\texttt{GT\_TOWER\_CHECK}&180&0&0\\
\texttt{LM7\_PUSH}&218&218&9\\
\texttt{LM8\_PUSH}&304&304&304\\
\hline
Total&1190&800&591\\
\hline
\end{tabular}
\end{center}
The nine LM7 blocks are E.1--E.7, E.58 and E.65. The $582$ blocks in
the complete A1, successor and LM8 main sources were available for
text screening. The targeted analytic passages were read; availability
does not mean that every mathematical proof was re-proved.

The flag file has $575$ records, representing $525$ distinct run/statement
pairs in $23$ runs. Its six-run subset has $166$ records and $152$
distinct pairs. The full audit table preserves all $800$ recovered rows,
and the curated table has $108$ entries, including the two other flagged
run examples, three intake items, the manuscript Artin passage and
Finding A06. Declared register dependencies give $201$ rows with a path
to one of $38$ curated qualitative prime-selection inputs. These paths
are recorded dependencies, not a newly certified proof graph.

For the two missing full registers, exactly $47$ A1-successor and $50$
GT identifiers are visible in the supplied summaries or flags. They
are listed separately; no unseen rows are invented. The remaining
coverage gap is therefore $390$ full register rows, together with $209$
LM7 proof blocks whose statements were recovered.

\paragraph{Finding R02 (the RH/GRH flags are false positives for conditional dependence).}
There are just two flags labelled RH/GRH. In LM8 A.69, the abbreviation
``DS/RH'' is resolved by the full proof as Deuring--Shafarevich and
Riemann--Hurwitz, not the Riemann Hypothesis. In
\texttt{S6\_SERIES\_GRIND}, row \texttt{maps:src:F.2.36}, the attached
context explicitly says that the class and unit groups were certified
with \texttt{GRH=false}. Its underlying computation was not recovered,
so this is an audit of the reported scope, not a new certification.
Neither flag establishes a formerly conditional theorem that is now
unconditional. Conversely, their absence from other rows does not show
analytic independence: the recovered LM8 proofs contain many
qualitative Chebotarev arguments.

\subsection{Qualitative prime choice and its quantitative refinements}

The finite-Galois qualitative Chebotarev theorem is unconditional:
a conjugacy class $C$ has density $|C|/|G|$. The exact source statement
checked here is Milne, \emph{Class Field Theory}, v4.03,
Chapter V, Theorem 3.23, printed p.~164.
The existing existence and density uses below therefore require neither
stipulated GRH hypothesis. A new least-prime version follows from
Finding G01 after the actual finite field and excluded primes have
been specified. Its exact analytic hypotheses are GRH for
\[
 \zeta_{\widetilde L}(s),\quad\zeta_{\mathbb Q}(s),\qquad
 \widetilde L=L\!\left(\sqrt{\operatorname{rad}
                     (2\prod_{p\in S}p)}\right).
\]
For a number-field base one uses the normal closure over $\mathbb Q$
and the rational primes below the excluded places. The stipulated
all-number-field Hecke hypothesis supplies these zeta functions as
trivial Hecke characters. Dirichlet GRH supplies this particular
argument when $L/\mathbb Q$ is abelian. No general Artin holomorphy
assumption is introduced.

\begin{longtable}{p{.22\linewidth}p{.43\linewidth}p{.26\linewidth}}
\hline
Rows & Actual finite field or predicate & Quantitative family\\
\hline
\endhead
LM8 A.16 & $M/\mathbb Q$ finite Galois containing the normalizer field;
complex-conjugation class. & $\zeta_M$, with G01 exclusions.\\[3pt]
LM8 B.20, B.22 & Normal closure of the specialized finite torsor fibre,
after the quotient, deck-constant field $K$, and Hilbert point $a$
are chosen. & Zeta of this normal closure.\\[3pt]
LM8 C.30; D.1; D.58 & Respectively
$\mathbb Q(\zeta_{12})$ with class $5$,
$\mathbb Q(\zeta_{5^r})$ with class $1$,
$\mathbb Q(\zeta_{30})$ with class $17$. & Dirichlet $L$-functions
modulo $12$, $5^r$, $30$, with finite exclusions.\\[3pt]
LM8 D.8 & $\mathbb Q(E_c[5],\mu_5)$; completely split primes.
& Zeta of this division field.\\[3pt]
LM8 D.11 & Kernel field of
$\rho_{M(-1)}\bmod\ell^{k+1}$ over the stated number field.
& Zeta of its $\mathbb Q$-normal closure.\\[3pt]
LM8 D.3, D.28 & Kernel field of
$g\mapsto(s_0(g)\bmod5,\chi_5(g)\bmod5^{r+1})$;
target $(\exp C,-1+5^r)$ in the actual Heisenberg image.
& Zeta of this actual image field.\\[3pt]
LM8 D.44, D.48, D.62, D.63, D.65 & Joint kernel field of the actual
finite Alexander action on $W_r$ and $\chi_5\bmod5^r$;
D.63 uses $r+1$. & Zeta of the joint image field.\\[3pt]
LM8 D.60, D.61, D.66 & Kernel field of the finite affine homomorphism
$g\mapsto(z(g),\rho_r(g))$ with additive coefficient module $W_r$.
& Zeta of the actual affine image field.\\[3pt]
LM8 D.67, D.69--D.71, D.73 & The same construction with finite
nonabelian coefficient group $H_r$ and its actual action;
arbitrary number field in D.70. & Zeta of the $\mathbb Q$-normal closure.\\[3pt]
LM8 D.78--D.87 & At each original finite Cartier-dual coefficient
module, the actual finite affine cocycle/action field.
& Zeta of that finite field.\\[3pt]
LM8 E.1, E.2 & $K=\mathbb Q(i)$ and $K(P)/K$ for a generator
$P$ of $E_6[12]$; degree $16$ in the source calculation.
& Finite-order Hecke $L_K(s,\chi)$, or $\zeta_{K(P)}$.\\[3pt]
LM7 A.16 & $\mathbb Q(\zeta_9,\sqrt[3]3)$;
$q\equiv4,7\pmod9$ and $3$ noncubic modulo $q$.
& Zeta of this degree-$18$ field.\\[3pt]
LM7 A.25 & $\mathbb Q(\zeta_9)$; classes $2,5$.
& Dirichlet $L(s,\chi)$ modulo $9$.\\
\hline
\end{longtable}

The LM7 proof sources for A.16 and A.25 are missing. Their fields are
reconstructed from the full register predicates, not from an asserted
recheck of their group presentations. Over $\mathbb Q(\zeta_3)$, the
extensions adjoining $\zeta_9$ and $\sqrt[3]3$ are independent cubic
extensions; four of the $18$ elements of the full Galois group satisfy
the conjugacy-stable A.16 predicate. This gives density $2/9$.
A.25's two classes among the six units modulo $9$ give density $1/3$.

For an additive or nonabelian cocycle, the finite homomorphism is the
actual graph $g\mapsto(z(g),\rho(g))$. A witness $\delta$ in the
action kernel with $z(\delta)\ne0$, respectively $\ne1$, gives a
nonempty conjugacy class in that actual image. One cannot replace it
by an arbitrarily chosen element of a larger formal affine group.
The recovered LM8 proofs make this distinction explicitly.

LM8 A.20 provides a separate logical warning: for a nontrivial quadratic
extension $L/K$, the inert primes have positive density and are
disjoint from the split primes, although the split set becomes
density one after pullback to $L$. This unconditional counterexample
does not refute the stronger arithmetic theorem for which a stronger
density hypothesis might suffice.

Stagewise effective prime choice does not bound a preceding Hilbert
specialization, a separating finite quotient, or the first obstruction
depth of an inverse system. LM7 A.43 and B.6 already give finite
field enumerations at each fixed depth from degree and discriminant
bounds; they expressly assert no uniform stopping rule for every
positive infinite-tower instance.

\subsection{Primitive roots, geometric existence and finite-field weights}

\paragraph{Finding R03 (fixed-prime criteria are not Artin-existence theorems).}
A1 B.3 and successor A.8, A.10, A.17 and A.35 use a hypothesis such as
$\operatorname{ord}_\ell(p)=\ell-1$ at a specified prime. This finite
condition makes $\Phi_\ell$ irreducible over $\mathbb F_p$ and its
nontrivial character values one Frobenius orbit. It does not produce
a prime satisfying the additional geometric or numerical conditions.
Successor A.12 supplies a genus-three counterexample to extending a
genus-two point-only test, even with $2$ primitive modulo $3$.

For the separate assertion of infinitely many primitive-root primes,
Hooley requires the family
\[
 K_n=\mathbb Q(\zeta_n,p^{1/n}),\qquad n\text{ squarefree},
 \qquad \zeta_{K_n}(s).
\]
The theorem was checked in \S7, printed p.~47, of the Russian
translation of his 1967 paper. The fields are generally nonabelian
over $\mathbb Q$, so the stipulated Dirichlet hypothesis does not
supply this family. The stipulated all-field Hecke hypothesis does.
Its arithmetic conclusion gives no theorem about intersection with a
curve-dependent theta-support set. The universal primitive-root theta
condition is in fact refuted by Finding A06. That example has
$p$-rank zero, and leaves the positive-$p$-rank restricted problem
unresolved. It is not a counterexample to A1.

The elementary density obstruction also remains useful: when
$\ell\equiv1\pmod8$, both $2$ and $-1$ are squares, so
$\langle2,-1\rangle$ is a proper subgroup of
$\mathbb F_\ell^\times$. A positive-density support set can therefore
be disjoint from the desired arithmetic condition. For torsion order
$\ell^a$, primitivity modulo $\ell$ also does not automatically
give primitivity modulo $\ell^a$.

A1 B.3 and successor A.11 instead use finite-field $L$-polynomials:
\[
 L(C,\chi,T)=\det(1-\operatorname{Frob}_qT
                  \mid H^1(C_{\bar k},\mathcal L_\chi)).
\]
Their nontrivial rank-one finite-order eigensheaves have $2g-2$
eigenvalues of absolute value $\sqrt q$. The resulting pair estimate is
\[
 \sum_bN_b(N_b-1)
 \le {N^2\over|A|}
   +{|A|-1\over|A|}(2g-2)^2q-N.
\]
Deligne's \emph{Weil I}, formula (1.5.1) and Theorem (1.6), are
unconditional finite-field inputs. Neither number-field GRH hypothesis
supplies the still missing quotient or prescribed-point existence in
A1 B.5 and successor A.15. Successor A.38 explicitly distinguishes
closed-point Chebotarev from prime-order theta avoidance on an entire
nonzero cyclic orbit.

The least-prime results A1 B.27, B.30 and B.34, and successor
A.19, A.24, A.31, A.72, A.74, A.91--A.93 and A.95, concern specified
curves or explicitly bounded finite censuses. They use finite
Frobenius/Cartier calculations and finite exclusions, not analytic
least-prime theorems. A1 B.34 is a \emph{field-first} minimum, not a
least prime over arbitrary extensions. A1 B.32 and successor A.96
retain unresolved cases. A1 B.4 decides a fixed prime over all
constant extensions; it does not bound the search over all primes.
Successor A.16 rules out a bound depending only on $(p,g)$.

\subsection{Other false positives and verification scope}

LM8 A.32--A.33 use Serre's qualitative non-CM open-image theorem;
their commutator Zariski-density conclusion was already unconditional.
The recovered proofs explicitly identify
$F_a/\mathbb Q:y^2=x^3-3x+a$, for $a=6,-6$, with $j=-216$;
the later choice of a number field $L$ is a restriction of this
representation. Thus there is no unresolved descent condition.
Finding G06 applies to these actual factors, for which
$\operatorname{rad}(2N_{F_a})=6$. Under $\mathrm H_M$, the mod-$\ell$
image over $\mathbb Q$ is $\operatorname{GL}_2(\mathbb F_\ell)$ for
$\ell>964\log6+5760$, and its restricted image over $L$ has index
at most $[L:\mathbb Q]$. Full image over $L$ is not asserted.
Independently of $\mathrm H_M$, the cited Lombardo theorem over
$\mathbb Q$ and the same restriction argument give
\[
 [\operatorname{GL}_2(\widehat{\mathbb Z}):\rho_{F_a,\infty}(G_L)]
 <[L:\mathbb Q]\,\exp(10^{21483})
       \max\{1,h_F(F_a)\}^{4.8\cdot10^{10}}.
\]
Here $h_F$ is the stable Faltings height in the cited theorem's
convention. These refinements do not change the original rows' status.

LM8 A.46--A.48 use Faltings's crystalline comparison; D.25 uses
Faltings's finiteness of rational points. Neither is a finite
Faltings--Serre trace-comparison test. The flagged SA07 F.3.6 excerpt
instead refers to a Hom theorem, and its full proof is missing.

The primitive roots in LM8 A.54 are generators of the two fixed
finite fields of orders $41$ and $61$; A.55 chooses a root of unity;
D.93 uses that $2$ generates $\mathbb Z_5^\times$. D.31 uses the
Artin braid action. The phrase ``vari\'et\'e \'el\'ementaire d'Artin''
in scan 037 is geometric, not a primitive-root conjecture.
GT C.1.21 takes the least \emph{index} with
$v_\ell((J+4)!)\ge N$. Its displayed definition is not a least-prime
problem, and its full proof remains unavailable.

The retained B.23 conductor enumeration and independent Bezout verifier
were rerun successfully. The latter recomputed the division-polynomial
compositions and the complete Bezout identities at conductors $191$
and $229$, after a missing runtime dependency was supplied. This is a
rerun of an existing independent implementation, not a newly written
one. A separate new elementary implementation independently checked
the six primes, multiplicative orders and all finite differences used
in A06. These events are labelled ``computation or formal certificate'';
targeted source reading is labelled ``used and rechecked in this run''
with the precise scope stated. All events are dated 2026-11-21 and name
GPT-6 Astra Pro; no unavailable effort setting is invented.

\paragraph{Unrecovered sources.}
Exact checkpoint and register names, timestamp variants, and broad
title searches were attempted for the A1-successor and GT returns.
No matching full source was recovered, and no nested matching archive
was present in the recovered bundles. The documented Personal Context
route returned ``Personal context is disabled''; it was blocked,
not a zero-result search. The remaining concrete route is to provide
the two original checkpoint archives or their full registers and papers,
or an exact identity for a differently titled upload.
\texttt{SOURCE\_RECOVERY\_MANIFEST.tsv} retains the exact recovered
identities, creation dates, sizes and SHA-256 hashes without signed
download URLs. The ledger's coverage statement must remain
$800/1190$ registers and $591/1190$ available source blocks until those
missing sources are supplied.
